
\documentclass[11pt]{article}

\usepackage{amsmath, amssymb, amsthm}
\usepackage[margin=1in]{geometry}


% Define the environment: \newtheorem{name}{Label}[reset_counter]
\newtheorem{theorem}{Theorem}[section] 
\newtheorem{lemma}[theorem]{Lemma} % Shares numbering with theorems
\newtheoremstyle{mydef}
{\topsep}{\topsep}%
{}{}%
{\bfseries}{}
{\newline}
{%
  \rule{\textwidth}{0.4pt}\\*%
  \thmname{#1}~\thmnumber{#2}\thmnote{\ -\ #3}.\\*[-1.5ex]%
  \rule{\textwidth}{0.4pt}}%

\theoremstyle{mydef}
\newtheorem{definition}{Definition}
\title{Calculus II Survival Sheet}
\author{}
\date{}

\begin{document}

\maketitle

\section*{Volumes of Revolution}

\subsection*{Disk Method}
\[
V=\pi\int_a^b R^2\,dx
\]

\subsection*{Washer Method}
\[
V=\pi\int_a^b (r_{outer}^2-r_{inner}^2)\,dx
\]

\subsection*{Shell Method}
\[
V=2\pi\int_a^b (\text{radius})(\text{height})\,dx
\]

Rule:
\begin{itemize}
\item Parallel slices $\rightarrow$ Washers
\item Perpendicular slices $\rightarrow$ Shells
\end{itemize}

\section*{Arc Length}

\[
L=\int_a^b \sqrt{1+(y')^2}\,dx
\]


\section*{Integration Techniques}

\subsection*{Integration by Parts}
\[
\int u\,dv=uv-\int v\,du
\]

\subsection*{Trigonometric Integrals}

Odd power of sine:
save one $\sin x$.

Odd power of cosine:
save one $\cos x$.

Even powers:
use identities
\[
\sin^2x=\frac{1-\cos2x}{2},\quad
\cos^2x=\frac{1+\cos2x}{2}
\]

\subsection*{Trig Substitution}

\[
\sqrt{a^2-x^2}\Rightarrow x=a\sin\theta
\]
\[
\sqrt{a^2+x^2}\Rightarrow x=a\tan\theta
\]
\[
\sqrt{x^2-a^2}\Rightarrow x=a\sec\theta
\]

Always draw a triangle to back-substitute.


\subsection*{Partial Fractions}

Used to integrate rational functions:
\[
\int \frac{P(x)}{Q(x)}\,dx
\]

\textbf{Requirement:}
\[
\deg(P)<\deg(Q)
\]
If not, perform polynomial long division first.

\bigskip
\textbf{Step 1: Factor the Denominator}

Factor $Q(x)$ completely over the reals.

\subsubsection*{Case 1: Distinct Linear Factors}

If
\[
Q(x)=(x-a)(x-b)
\]

then
\[
\frac{P(x)}{(x-a)(x-b)}
=\frac{A}{x-a}+\frac{B}{x-b}
\]

Example:
\[
\frac{3x+5}{(x-1)(x+2)}
=\frac{A}{x-1}+\frac{B}{x+2}
\]

\subsubsection*{Case 2: Repeated Linear Factors}

If
\[
Q(x)=(x-a)^n
\]

include every power:
\[
\frac{P(x)}{(x-a)^n}
=\frac{A_1}{x-a}
+\frac{A_2}{(x-a)^2}
+\cdots
+\frac{A_n}{(x-a)^n}
\]

Example:
\[
\frac{2x+1}{(x-3)^3}
=
\frac{A}{x-3}
+\frac{B}{(x-3)^2}
+\frac{C}{(x-3)^3}
\]

\subsubsection*{Case 3: Irreducible Quadratic Factors}

If the quadratic cannot factor over the reals:
\[
x^2+bx+c
\]

use a linear numerator:
\[
\frac{Ax+B}{x^2+bx+c}
\]

Example:
\[
\frac{x}{(x-1)(x^2+4)}
=
\frac{A}{x-1}
+\frac{Bx+C}{x^2+4}
\]

\subsubsection*{Case 4: Repeated Irreducible Quadratics}

If
\[
(x^2+1)^2
\]

then
\[
\frac{Ax+B}{x^2+1}
+\frac{Cx+D}{(x^2+1)^2}
\]

\subsection*{Improper Integrals}
Improper integrals are defined using limits.

%------------------------------------------------

\subsubsection*{Type I: Infinite Interval}

\textbf{Infinite upper bound}
\[
\int_a^{\infty} f(x)\,dx
=\lim_{b\to\infty}\int_a^b f(x)\,dx
\]

\textbf{Infinite lower bound}
\[
\int_{-\infty}^{a} f(x)\,dx
=\lim_{b\to-\infty}\int_b^a f(x)\,dx
\]

\textbf{Both directions}
\[
\int_{-\infty}^{\infty} f(x)\,dx
=
\int_{-\infty}^{c}f(x)\,dx
+
\int_c^{\infty}f(x)\,dx
\]

Each limit must converge separately.

\subsubsection*{Type II: Infinite Integrand (Vertical Asymptote)}

If $f(x)$ is discontinuous or undefined at an endpoint or inside the interval, rewrite using limits.

\textbf{At right endpoint $x=b$:}
\[
\int_a^b f(x)\,dx
=
\lim_{t\to b^-}\int_a^t f(x)\,dx
\]

\textbf{At left endpoint $x=a$:}
\[
\int_a^b f(x)\,dx
=
\lim_{t\to a^+}\int_t^b f(x)\,dx
\]


\textbf{If discontinuous at $x=c$ inside $(a,b)$:}
\[
\int_a^b f(x)\,dx
=
\lim_{t\to c^-}\int_a^t f(x)\,dx
+
\lim_{t\to c^+}\int_t^b f(x)\,dx
\]

Both limits must converge for the integral to converge.
\section*{Differential Equations (Separable + Initial Value Problems)}

A differential equation relates a function and its derivatives.

We focus on \textbf{separable equations}:
\[
\frac{dy}{dx} = g(x)h(y)
\]



\subsubsection*{Step 1: Separate Variables}

Rearrange to isolate $x$ and $y$:

\[
\frac{1}{h(y)}\,dy = g(x)\,dx
\]

This is the key move: get all $y$ terms with $dy$, all $x$ terms with $dx$.



\subsubsection*{Step 2: Integrate Both Sides}

\[
\int \frac{1}{h(y)}\,dy = \int g(x)\,dx
\]

Do NOT forget the constant:
\[
+ C
\]

Usually written on one side after integrating (two constants become one new constant).


\subsubsection*{Step 3: Solve for $y$ (if possible)}

After integrating, isolate $y$ explicitly if you can.

Sometimes you leave it implicit if solving is messy.



\subsubsection*{Initial Value Problems (IVPs)}

Given:
\[
\frac{dy}{dx} = g(x)h(y), \quad y(x_0)=y_0
\]

Steps:
\begin{enumerate}
\item Separate variables
\item Integrate both sides
\item Apply initial condition $(x_0, y_0)$
\item Solve for constant $C$
\item Solve for $y$ (if possible)
\end{enumerate}

\section*{L'Hôpital's Rule}

Used to evaluate limits of indeterminate forms:

\[
\frac{0}{0} \quad \text{or} \quad \frac{\infty}{\infty}
\]

If
\[
\lim_{x\to a} \frac{f(x)}{g(x)} \text{ is } \frac{0}{0} \text{ or } \frac{\infty}{\infty},
\]

then:
\[
\lim_{x\to a} \frac{f(x)}{g(x)}
=
\lim_{x\to a} \frac{f'(x)}{g'(x)}
\]

Apply repeatedly if needed.

\textbf{Important:} Only use if the original limit is truly indeterminate.

\subsubsection*{Other Indeterminate Forms (must rewrite first)}

L'Hôpital's Rule does NOT directly apply to:

\[
0 \cdot \infty,\quad \infty - \infty,\quad 0^0,\quad 1^\infty,\quad \infty^0
\]

These must be algebraically rewritten first.



\subsubsection*{Indeterminate Powers Strategy}

For expressions of the form:
\[
f(x)^{g(x)}
\]

Take natural logs:

\[
y = f(x)^{g(x)}
\quad \Rightarrow \quad
\ln y = g(x)\ln(f(x))
\]

Then evaluate:
\[
\lim \ln y = \lim g(x)\ln(f(x))
\]

Now apply algebra + L'Hôpital if needed.

Finally exponentiate:

\[
\lim y = e^{\lim \ln y}
\]

\section*{Sequences}

\begin{definition}
A \textit{sequence} is a function $a_n : \mathbb{N} \to \mathbb{R}$.

\[
\{a_n\} = a_1, a_2, a_3, \dots
\]
This is an infinitely long ordered list. 
\end{definition}
\subsection*{Limit of a Sequence}

\begin{definition}
    The sequence $\{a_n\}$ has the \textit{limit} L and we write
\[
\lim_{n \to \infty} a_n = L.
\]
If $lim_{n \to \infty} a_n$ exists, we say the sequence \textit{converges}. Otherwise, we say the sequence \textit{diverges}.
\end{definition}

\begin{theorem}
    If 
\[
\lim_{n \to \infty} f(x) = L.
\] and $f(n)=a_n$ where n is an integer, then $
\lim_{n \to \infty} a_n = L.$
\end{theorem}

\textbf{Example:} $1,-1,1,-1,...$ is of the form $(-1)^{n+1}$\\
\textbf{Example:} $-1,1,-1,1,...$ is of the form $(-1)^n$
\subsection*{Algebraic Limit Laws}
If $\lim a_n = A$ and $\lim b_n = B$, then:
\begin{align*}
$\lim (a_n + b_n) &= A+B$ \\
$\lim (a_n b_n) &= AB$ \\
$\lim \frac{a_n}{b_n} &= \frac{A}{B}, \quad B\neq0$
\end{align*}

\subsection*{Squeeze Theorem}
If
\[
a_n \le b_n \le c_n
\]
and
\[
\lim a_n = \lim c_n = L,
\]
then
\[
\lim b_n = L.
\]

\section*{Infinite Series}

\begin{definition}
    An infinite series is
\[
\sum_{n=1}^{\infty} a_n.
\]
\end{definition}

\begin{definition}
Define partial sums:
\[
S_N=\sum_{n=1}^N a_n.
\]
\end{definition}

\textbf{Example:}
These are partial sums:\\
$s_1=a_1$\\
$s_2=a_1+a_2$\\
In general we have:\\
$s_n=a_1+a_2+...+a_n$\\

\noindent The series converges if $\lim_{N\to\infty} S_N$ exists.

\begin{definition}
    Given a series $\sum_{n=1}^\infty a_n=a_1+a_2+a_3+...,$ let $s_n$ denote its $n$th partial sum.:\\
\[
s_n=\sum_{i=1}^n a_i=a_1+a_2+a_3+...+a_n.
\]
If the seqeuence $\{s_n\}$ is convergent and $\lim_{n\to\infty} s_n=s$ exists as a real number, then the series  $\sum_{n=1}^\infty a_n$ is called \textit{convergent} and we write:\\
\[
a_1+a_2+a_3+...+a_n=s.
\]
or
\[
\sum_{i=1}^\infty a_i = s
\]
The number $s$ is called the \textit{sum} of the series. If the sequence $\{s_n\}$ is divergent, the the series is called \textit{divergent}.
\end{definition}
\subsection*{Divergence Test}
If
\[
\lim_{n\to\infty} a_n \neq 0,
\]
then the series diverges. If it is 0, no conclusion can be drawn.\\
Note: This test can be used on any series.

\textbf{Linearity of Series}
\[
\sum_{n=1}^{\infty} c \, a_n &= c \sum_{n=1}^{\infty} a_n, \quad c \in \mathbb{R}
\]
\[\sum_{n=1}^{\infty} (a_n \pm b_n) &= \sum_{n=1}^{\infty} a_n \pm \sum_{n=1}^{\infty} b_n
\]
\section*{Telescoping Series}
A telescoping series has cancellation between terms:
\[
\sum (a_n - a_{n+1})
\]
Steps:
\begin{itemize}
\item Rewrite using partial fractions if needed
\item Expand partial sums
\item Cancel terms
\item Take limit of remaining boundary terms
\end{itemize}
\section*{Geometric Series}
A \textit{geometric series} is of the form:
\[
\sum_{n=1}^{\infty} ar^{n-1}=a+ar+ar^2+ar^3+...+ar^{n-1}+...
\]
where $a\neq 0$\\
Each term is obtained from the preceding one by multiplyig it by the \textit{common ratio (r)}.\\
\\
The geometric series 
\[
\sum_{n=1}^{\infty} ar^{n-1}=a+ar+ar^2+ar^3+...+ar^{n-1}+...
\]
is convergent if $|r|<1$ and its sum is:
\[
\frac{a}{1-r}.
\]
If $|r|\ge1$ the geometric series is divergent.\\
Note: This is specifically for the geometric series.

\section*{p-Series}
A \textit{p-series} is of the form:
\[
\sum_{n=1}^{\infty} \frac{1}{n^p}
\]

\begin{itemize}
\item Converges if $p>1$
\item Diverges if $p\le1$
\end{itemize}

\textbf{Harmonic Series:}
\[
\sum_{n=1}^{\infty} \frac{1}{n} = 1 + \frac{1}{2} + \frac{1}{3} + \frac{1}{4} + \cdots
\]

Using the p-series test: $\sum_{n=1}^{\infty} \frac{1}{n^p}$ converges if $p>1$, diverges if p \le 1.

\text{Here, } p = 1, \text{ so the series diverges.}

\section*{Integral Test}

Let $f$ be continuous from $[1,\infty)$, positive, decreasing and $a_n=f(n)$.

Then
\[
\sum a_n \text{  converges if and only if } \int_1^\infty f(x)\,dx \text{ converges}
\]

\[
\sum a_n \text{  diverges if and only if } \int_1^\infty f(x)\,dx \text{ diverges}
\]
either both converge or both diverge.

\section*{Comparison Tests}

\subsection*{Direct Comparison Test}
Suppose that $\sum a_n$ and $\sum b_n$ are series with positive terms.

\begin{itemize}
\item If $\sum b_n$ converges and $a_n\leq b_n  \forall n$ $\Rightarrow \sum a_n$ converges.
\item If $\sum b_n$ diverges and $a_n\geq b_n \forall n$ $\Rightarrow \sum a_n$ diverges.
\end{itemize}

\subsection*{Limit Comparison Test}
Suppose  that $\sum a_n$ and $\sum b_n$ are series with postivie terms. If 
\[
\lim_{n\to\infty}\frac{a_n}{b_n}=c,
\quad 0<c<\infty,
\]
then either both series converge or both diverge.

\section*{Alternating Series}

\begin{definition}
    An \textit{alternating series} is a series whose terms are alternately positive and negative.
\end{definition}

\subsection*{Alternating Series Test}
For
\[
\sum (-1)^n b_n,
\]
if:
\begin{itemize}
\item $b_{n+1}\leq b_n \forall n$,
\item $\lim_{n \to \infty} b_n = 0$,
\end{itemize}
then the series converges.\\
\textbf{Note:} $b_n$ is the absolute value of $a_n$


\section*{Absolute and Conditional Convergence}

\begin{definition}
    A series is \textit{absolutely convergent} if $\sum |a_n|$ converges.
\end{definition}
\begin{definition}
    A series is \textit{conditionally convergent} if it is convergent but not absolutely convergent ($\sum |a_n|$ diverges).
\end{definition}

\begin{theorem}
If a series $\sum a_n$ is absolutely convergent, then it is convergent.
\end{theorem}


\section*{Ratio Test}

Let
\[
L=\lim_{n\to\infty}\left|\frac{a_{n+1}}{a_n}\right|.
\]

\begin{itemize}
\item If $\lim_{n\to\infty}\left|\frac{a_{n+1}}{a_n}\right|=L<1$ the series $\sum a_n$ is absolutely convergent (So it is convergent)
\item If $\lim_{n\to\infty}\left|\frac{a_{n+1}}{a_n}\right|=L>1$ the series $\sum a_n$ is divergent
\item If $\lim_{n\to\infty}\left|\frac{a_{n+1}}{a_n}\right|=L=1$ then the ratio test is inconclusive for the series $\sum a_n$.
\end{itemize}

\section*{Power Series}

\[
\sum_{n=0}^{\infty} c_n(x-a)^n
\]
\subsection*{Radius of Convergence}
Find $R$ using Ratio Test.

\begin{itemize}
\item Converges for $|x-a|<R$
\item Diverges for $|x-a|>R$
\item Test endpoints separately
\end{itemize}
\subsection*{Power Series Operations}
You can differentiate and integrate term-by-term:
\begin{align*}
\frac{d}{dx}\sum c_n(x-a)^n &= \sum n c_n (x-a)^{n-1} \\
\int \sum c_n(x-a)^n dx &= \sum \frac{c_n}{n+1}(x-a)^{n+1}
\end{align*}

\subsection*{Interval of Convergence Notes}
Always express final answers as:
\begin{itemize}
\item $(a-R, a+R)$
\item $[a-R, a+R]$
\item mixed endpoints like $(a-R, a+R]$
\end{itemize}
\section*{Taylor Series}

\[
f(x)=\sum_{n=0}^{\infty}
\frac{f^{(n)}(a)}{n!}(x-a)^n
\]

\subsection*{Maclaurin Series (Taylor Series with $a=0$)}

\begin{align*}
e^x &= \sum_{n=0}^{\infty}\frac{x^n}{n!}, \quad R=\infty \\
\sin x &= \sum_{n=0}^{\infty}(-1)^n\frac{x^{2n+1}}{(2n+1)!}, \quad R=\infty \\
\cos x &= \sum_{n=0}^{\infty}(-1)^n\frac{x^{2n}}{(2n)!}, \quad R=\infty \\
\frac{1}{1-x} &= \sum_{n=0}^{\infty}x^n, \quad |x|<1,\; R=1 \\
\ln(1+x) &= \sum_{n=1}^{\infty}(-1)^{n+1}\frac{x^n}{n}, \quad |x|<1,\; R=1 \\
(1+x)^k &= \sum_{n=0}^{\infty} \binom{k}{n} x^n, \quad |x|<1,\; R=1 \\
\tan^{-1}x &= \sum_{n=0}^{\infty}(-1)^n \frac{x^{2n+1}}{2n+1}, \quad |x|\le 1,\; R=1
\end{align*}

\section*{Parametric Functions}

A parametric curve is defined by:
\[
x = x(t), \quad y = y(t)
\]

The parameter $t$ traces the curve.

\subsubsection*{Slope of a Parametric Curve}

\[
\frac{dy}{dx} = \frac{\frac{dy}{dt}}{\frac{dx}{dt}}
\]

\textbf{Important:}
\[
\frac{dx}{dt} \neq 0
\]

\subsubsection*{Tangent Line (Parametric)}

At $t=t_0$:

\begin{align*}
x_0 &= x(t_0) \\
y_0 &= y(t_0)
\end{align*}

Slope:
\[
m = \frac{dy/dt}{dx/dt}\Bigg|_{t=t_0}
\]

Equation:
\[
y - y_0 = m(x - x_0)
\]



\subsubsection*{Area (Parametric)}

\[
A = \int_{t=a}^{t=b} y(t)\frac{dx}{dt}\,dt
\]


\subsubsection*{Arc Length (Parametric)}

\[
L = \int_a^b \sqrt{\left(\frac{dx}{dt}\right)^2 + \left(\frac{dy}{dt}\right)^2}\,dt
\]


\section*{Polar Functions and Regions}

Polar form:
\[
r = f(\theta)
\]

Conversion:
\[
x = r\cos\theta, \quad y = r\sin\theta
\]


\subsubsection*{Slope in Polar Coordinates}

\[
\frac{dy}{dx}
=
\frac{\frac{dr}{d\theta}\sin\theta + r\cos\theta}
{\frac{dr}{d\theta}\cos\theta - r\sin\theta}
\]


\subsubsection*{Area in Polar Coordinates}

\[
A = \frac{1}{2}\int_{\alpha}^{\beta} r^2\,d\theta
\]

\textbf{Between two curves:}
\[
A = \frac{1}{2}\int_{\alpha}^{\beta} \left(r_{\text{outer}}^2 - r_{\text{inner}}^2\right)\,d\theta
\]


\subsubsection*{Arc Length (Polar)}

\[
L = \int_{\alpha}^{\beta} \sqrt{r^2 + \left(\frac{dr}{d\theta}\right)^2}\,d\theta
\]



\subsubsection*{Finding Area Between Polar Curves}

Steps:
\begin{enumerate}
\item Find intersection angles ($r_1 = r_2$)
\item Determine outer and inner curve
\item use the formula
\end{enumerate}
\end{document}
